\documentclass[sigconf,screen]{acmart}

\begin{document}

\title{🤖 rocq-robot: Neurosymbolic Theorem Proving via LLM-LSP Integration}

\author{Gavin Mendel-Gleason}
\affiliation{%
  \institution{Scidonia}
  \city{Dublin}
  \country{Ireland}
}
\email{gavin@scidonia.com}

\maketitle

\section*{Demo Abstract}

We demonstrate \textbf{rocq-robot}, a Model Context Protocol (MCP) server enabling seamless integration between Large Language Models (LLMs) and the Rocq proof assistant (formerly Coq) through the Language Server Protocol (LSP). This tool embodies \emph{neurosymbolic programming}: real-time collaboration between neural networks (LLMs providing heuristic guidance) and symbolic reasoning systems (proof assistants ensuring correctness).

\paragraph{What We Will Demonstrate}

The demo will showcase:

\begin{enumerate}
\item \textbf{Interactive proof construction}: Live LLM-driven proof development with rocq-robot tools, showing how an AI agent navigates proof state, queries available lemmas, and applies tactics with immediate verification feedback.

\item \textbf{Speculative execution}: Demonstrating how LLMs can "try before committing"—testing tactics speculatively without file modifications, enabling efficient exploration of proof strategies.

\item \textbf{Autonomous theorem proving}: Presenting our completed type preservation proof for PCF with references (21 inductive cases, 7 helper lemmas)—entirely constructed by an LLM using only tool-based interactions, no human intervention.

\item \textbf{Error recovery}: Showing how failed tactics provide actionable feedback, allowing LLMs to adjust strategies and retry—demonstrating the robustness of the neurosymbolic approach.

\item \textbf{Project-aware development}: Illustrating automatic workspace switching when working across multiple Rocq projects with different dependencies.
\end{enumerate}

\paragraph{Technical Innovation}

rocq-robot provides 20+ MCP tools organized into:
\begin{itemize}
\item \textbf{Proof interaction}: Focus on proofs, query goals, insert tactics with automatic bullet management, undo operations
\item \textbf{Knowledge search}: Find lemmas, check types, locate definitions, speculatively import libraries
\item \textbf{State management}: Low-level Pétanque API for fine-grained proof state control
\end{itemize}

Unlike existing LLM+theorem prover integrations (e.g., LeanDojo, Baldur), rocq-robot focuses on the \emph{interaction protocol} rather than model training, making it usable with any LLM through standardized tool interfaces.

\paragraph{Practical Value}

rocq-robot demonstrates that formal verification—traditionally requiring expert knowledge—can be made accessible through AI assistance while maintaining absolute correctness guarantees. Every proof step is verified by Rocq's kernel; the LLM provides creative guidance, the proof assistant ensures soundness.

\paragraph{Availability}

The tool is open-source (MIT license), production-ready, and actively used:
\begin{itemize}
\item GitHub: \url{https://github.com/scidonia/rocq-robot}
\item Documentation: Comprehensive guides, API reference, working examples
\item Requirements: Node.js 20+, Rocq/Coq 8.19+, rocq-lsp 0.2.0+
\end{itemize}

\paragraph{Demo Format}

We will provide:
\begin{itemize}
\item Live coding session showing LLM constructing proofs in real-time
\item Interactive Q\&A where attendees can suggest theorems for the LLM to prove
\item Technical deep-dive into the LSP↔MCP translation architecture
\item Discussion of future directions: multi-agent orchestration, cross-assistant support, proof repair
\end{itemize}

Attendees will leave understanding how neurosymbolic programming enables autonomous formal verification with guaranteed correctness.

\end{document}
